\section*{逻辑的形式化。}
\phantomsection
\addcontentsline{toc}{section}{逻辑的形式化。}

从我们掌握的有关公元前 5 世纪希腊哲学思想的（非常零散的）文本来看，总的印象是：这一时期被一种日益自觉的努力所支配，即把当时修辞学和数学如此成功地付诸实施的引导讨论的程序，推广到人类思想的全部领域——换句话说，去创立最一般形式的逻辑。哲学著作的语气此时发生了突变：在公元前 7 世纪和公元前 6 世纪，哲学家们只是断言或预言（至多依据同样含糊的类比勾画出含糊的论证），而从巴门尼德开始，尤其是到了芝诺，他们展开论辩，并力图提炼出可以作为其辩证法基础的一般原理：正是在巴门尼德那里可以找到排中律的最早陈述，而埃利亚的芝诺的那些“归谬法”证明则久负盛名。但芝诺写作于公元前 5 世纪中叶；而且，无论我们的文献有多少不确定之处，\(^{2}\) 很可能在这个时代，数学家们在他们自己的领域里已经在日常使用这些原理了。

如前文所说，追述这种逻辑在孕育过程中步步皆是的大量困难，以及由此产生的从埃利亚学派经智者派到柏拉图和亚里士多德的种种论争，并非我们的任务；我们在这里只拣出演讲艺术的勤勉耕耘及其后果之一——语言分析——在这一演变中所起的作用，这些发展通常被认为主要应归功于公元前 5 世纪的智者派。另一方面，数学的影响即使不总是被明确承认，也依然是显而易见的，在柏拉图和亚里士多德的著作中尤其如此。可以说，柏拉图几乎被数学迷住了；虽然他本人在这个领域并非创造者，但在他一生中的某个特定时期之后，他一直关注同时代数学家（其中许多人是他的朋友或学生）的发现，从此再未停止过一种最为直接的兴趣，甚至到了提出新研究方向的地步；因此，在他的著作中，数学不断地充当例证或模型（有时甚至像在毕达哥拉斯学派那里一样，滋养了他走向神秘主义的倾向）。至于他的学生亚里士多德，他不可能没有接受学园的学生所必需的最低限度的数学基础，人们甚至出版过一部由他著作中论及数学或涉及数学的文句辑成的集子 {[}153 d{]}；但他似乎从未下过大力气去了解他那个时代的数学动态，他在这个领域所引述的，只是很早以前就已普及的结果。这种错位在后世大多数哲学家那里只会愈加明显，他们当中许多人由于缺乏技术上的准备，真心诚意地自以为在内行地谈论数学，而实际上他们所指涉的不过是数学演进中早已被超越的一个阶段。

就逻辑而言，这一时期的最终产物是亚里士多德的宏篇巨制 {[}6{]}，其巨大功绩在于第一次成功地系统化并整理了在其前辈那里一直处于含糊或未定形的推理程序。\(^{3}\) 在这里，我们必须首先记住这部著作的总论点，即：一切正确的推理都可以化归为对一小批固定不变的规则的系统应用，而这些规则独立于所论对象的特殊本性（用字母来表示概念或命题——这大概是亚里士多德从数学家那里借用来的——清楚地表明了这种独立性）。但亚里士多德的注意力几乎完全集中在一种特殊类型的关系和逻辑连锁上，它们构成他所称的“三段论”：这主要是指一些若用集合论语言表述、今天会写成 \(A \subset B\) 或 \(A \cap B \neq \emptyset\) 形式的关系，\(^{4}\) 以及把这些关系或其否定借助如下格式串联起来的方式

\[
(A \subset B \text {   and   } B \subset C) \Rightarrow (A \subset C).
\]

亚里士多德对他那个时代的数学仍了解得足够多，因而看得出这类格式并不足以涵盖数学家所使用的全部逻辑运算，更不足以涵盖逻辑的其他应用（{[}6{]}，An. Pr., I, 46）。\(^{5}\) 至少，他对各种形式的“三段论”所作的深入研究（这一研究几乎全部用于澄清由推理所及对象的含糊或晦涩不断引起的种种困难）也使他有机会表述出求得一个命题的否定的规则（{[}6{]}，An. Pr., I, 46）。同样应归功于亚里士多德的，是他极精确地区分了“全称”命题与“特称”命题各自的作用，这是量词的最初雏形。\(^{6}\) 但是，他的著作（常被狭隘而缺乏理解地加以解释）直到 19 世纪很久以后仍有巨大影响，这种影响如何助长了哲学家对数学研究的忽视，并阻碍了形式逻辑的进步，这已是尽人皆知的事了。\(^{7}\)

然而，形式逻辑在古代仍继续取得进步，那是在作为逍遥学派对手的麦加拉学派和斯多葛学派的周围。遗憾的是，我们关于这些学说的信息全部来自二手，而且往往出自论敌或平庸注释者之口。这些逻辑学家取得的主要进展，看来在于建立了今天所说的“命题演算”：他们不像亚里士多德那样局限于 \(A \subset B\) 这种特殊形式的命题，而是对完全不确定的命题陈述规则。此外，他们对这些规则之间的逻辑联系分析得如此深入，以致知道如何从被立为“不可证明”的五条规则出发，用与现代方法非常类似的程序推出其余全部规则 {[}23{]}。可惜他们的影响相当短暂，他们的成果遂归于湮没，直到 19 世纪的逻辑学家重新发现它们为止。直到 17 世纪，亚里士多德始终是逻辑学界无可争议的宗师；众所周知，经院哲学家完全处于他的支配之下，而尽管他们对形式逻辑的贡献绝非微不足道，{[}25{]} 其中却不包含相对于古代哲学家成就而言在最高层次上的任何进步。

不过，在这里指出下面一点是恰当的：亚里士多德及其后继者的著作似乎并未对数学产生显著影响。希腊数学家沿着毕达哥拉斯学派及公元前 4 世纪后继者（泰奥多鲁斯、泰阿泰德、欧多克索斯）开辟的道路从事研究，而在表述其结果时显然并不拘泥于形式逻辑：只要把当时数学推理已经获得的灵活性和精确性，同亚里士多德逻辑尚处于非常初等的状况作一比较，这一事实就几乎不会令人惊讶。而当逻辑超越这一阶段时，指导其演变的又将是来自数学的新收获。

随着代数的发展，人们确实不可能不被形式逻辑规则与代数规则之间的类似所打动：二者共同的性质在于，都适用于（命题或数这样一些）并未精确确定的对象。而当 17 世纪代数记号在韦达和笛卡尔手中取得定型时，几乎立刻就出现了各种旨在表示逻辑运算的符号记号的尝试；但在莱布尼茨之前，这些尝试性的努力，例如埃里戈纳（1644 年）想用来记写初等几何证明的记号，或佩尔（1659 年）想用来记写算术证明的记号，都仍停留在非常表面的层次，没有带来数学推理分析上的任何进步。

到了莱布尼茨，我们面对的是一位同时也是第一流数学家的哲学家，他知道如何从自己的数学经验中提炼出使形式逻辑走出经院死胡同的思想萌芽。\(^{8}\) 作为一位举世罕见的博学通才、一个取之不竭的独创而多产的思想之源，莱布尼茨对逻辑的兴趣尤为浓厚，因为逻辑正处于他为使语言和思想形式化而设的宏伟计划的中心，而他毕生都在为这一计划不懈地工作。他自幼受经院逻辑的熏陶，却被这样一个（可上溯到雷蒙·吕勒的）想法所吸引：这种方法把一切人类概念还原为原始概念，构成一个“人类思想的字母表”，并以近乎机械的方式重新组合它们，从而得到一切真命题（{[}198 b{]}，第 VII 卷，第~185~页；参见 {[}69a{]}，第 II 章）。还在非常年轻的时候，他就形成了另一个独创得多的思想，即符号记号可以作为思想的“阿里阿德涅之线”而大有用处：\(^{9}\) “真正的 方法”，他说，“应当给我们一条阿里阿德涅之线，也就是说，一种粗略而有形的引导心灵的手段，就像几何中所画的线以及规定给算术学徒的那类运算一样。没有它，我们的心灵就无法沿着长路前行而不迷失。”（{[}198 b{]}，第 VII 卷，第~22~页；参见 {[}69 a{]}，第~90~页）。直到 25 岁上下，他对当时的数学都所知甚少，因此起初是以“通用语言”的形式提出他的计划（{[}69 a{]}，第 III 章）；但一与代数接触，他便把它采作其“普遍文字”的模型。他指的是这样一类符号语言：它能毫无歧义地表达一切人类思想，增强我们的推理能力，通过一种纯然机械的专注努力来避免错误，并且构造得使“连提出空想的人自己都不了解的那种空想，也不可能用这些字符写下来”（{[}198 a{]}，第 I 卷，第~187~页）。在莱布尼茨著作中论及这一宏伟计划及其实现将带来何种进步的无数段落里（参见 {[}69 a{]}，第 IV 章和第 VI 章），可以看到他对形式化语言的概念构想得何等清晰：一种纯粹的符号组合，其中重要的只是符号的连接，\(^{10}\) 以至一台机器就能产生出所有的定理，\(^{11}\) 而一切争论都可以通过简单的计算来解决（{[}198 b{]}，第 VII 卷，第~198-203~页）。即使这些希望今天看来有些过高，莱布尼茨数学工作的相当一部分——从他在无穷小演算的符号体系方面的工作开始（见第~190~页以下）——确实应当系于他思想中这种一以贯之的倾向，这一点仍然是真的；他自己对此有充分的自觉，还明确地把关于指数记号和行列式的想法（{[}198 a{]}，第 II 卷，第~204~页；参见 {[}69 a{]}，第~481-487~页）以及“几何演算”的概要（见第~50~页和第~61~页以下；参见 {[}69 a{]}，第 IX 章）也同他的“文字”联系在一起。但在他心中，本质的部分必须是符号逻辑，或者如他所说，是一个“Calculus ratiocinator”（推理演算），虽然他没有能创立这种演算，但我们至少看到他尝试了不下三次。第一次尝试时，他想到给每个“原始”词项关联一个素数，由若干原始词项组成的词项则用相应素数的乘积来表示；\(^{12}\) 他试图把通常的三段论规则翻译到这个系统中去，但遇到了由否定引起的重大复杂化（他很自然地想用变号来表示否定），于是很快放弃了这条道路（{[}198 c{]}，第~42-96~页；参见 {[}69 a{]}，第~326-344~页）。在后来的尝试中，他力图赋予亚里士多德逻辑一种更为代数化的形式；他有时保留记号 \(AB\) 来表示两个概念的合取，有时采用记号 \(A + B\)；\(^{13}\) 他（用乘法记号）写下了幂等律 \(AA = A\)，指出可以把命题“所有 A 都是 B”换成等式 \(A = AB\)，并且从这一点出发，用纯代数的计算就能重新得到亚里士多德的大部分规则（{[}198 c{]}，第~229-237~页和第~356-399~页；参见 {[}69 a{]}，第~345-364~页）；他还有了空集（“non Ens”）的想法，并且例如认识到了命题“所有 A 都是 B”与“A.(非 B) 不存在”的等价（出处同上）。此外，他注意到他的逻辑演算不仅适用于概念的逻辑，而且适用于命题的逻辑（{[}198 c{]}，第~377~页）。这样看来，他已经非常接近“布尔演算”。可惜的是，他似乎未能完全摆脱经院的影响；他不仅把自己演算的几乎唯一目标定为用他的记号去转录三段论的规则，\(^{14}\) 甚至为了完全恢复亚里士多德的规则而不惜牺牲自己最幸运的想法，哪怕这些规则中有些与空集观念不相容。\(^{15}\)

莱布尼茨的著作大部分直到 20 世纪初仍未发表，直接的影响因而甚微。整个 18 世纪以及 19 世纪初，有好几位作者（德·塞格纳、J. 兰伯特、普卢克凯、霍兰德、德·卡斯蒂隆、热尔戈纳）做了与莱布尼茨类似的尝试，但从未实质性地超出他停步的地方；这些工作收效甚微，以致他们中的大多数人对自己先行者的成果一无所知。\(^{16}\) 无论如何，G. 布尔——现代符号逻辑真正的创立者非他莫属 {[}29{]}——正是在同样的条件下从事写作的。他的关键想法在于系统地采取考虑“外延”的立场，即直接用集合来计算，把两个集合的交记作 xy，把 x 与 y 没有公共元素时二者的并记作 \(x + y\)。他还引入了一个记作 1 的“全集”（全体元素的集合）和记作 0 的空集，并把 x 的余集写作 1 - x。像莱布尼茨所做过的那样，他用关系 xy = x 来解释包含关系（他由此毫不费力地提取出对经典三段论规则的论证），而他对并与余集的记法也给他的系统带来了先行者们所缺乏的灵活性。\(^{17}\) 此外，通过把每个命题同它成立的那些“情形”的集合相关联，他把蕴涵关系解释为一种包含关系，于是他的集合演算便以这种方式为他提供了“命题演算”的那些规则。

19 世纪下半叶，布尔的系统成了一个活跃的逻辑学家学派的工作基础，他们在好几点上改进并完善了它。例如，杰文斯（1864 年）扩大了并运算 \(x + y\) 的含义，把它推广到 \(x\) 与 \(y\) 任意的情形；A. 德·摩根在 1858 年、C. S. 皮尔士在 1867 年证明了对偶律

\[
(\mathbf {C} A) \cap (\mathbf {C} B) = \mathbf {C} (A \cup B), (\mathbf {C} A) \cup (\mathbf {C} B) = \mathbf {C} (A \cap B); ^{18}
\]

德·摩根还在 1860 年着手研究关系，定义了二元关系的逆与复合（也就是与图上的运算 \(\overline{G}^{-1}\) 和 \(G_{1} \circ G_{2}\) 相对应的那些运算）。\(^{19}\) 所有这些工作都在施罗德卷帙浩繁而多产的巨著 {[}277{]} 中得到系统的阐述和发展。但颇为有趣的是，上面谈到的这些逻辑学家似乎对把自己的成果应用于数学并不很感兴趣；相反，布尔尤其是施罗德似乎以模仿经典代数的方法和问题（往往相当人为）来发展“布尔”代数作为自己的主要目标。这种态度的原因，无疑应当从如下事实中去寻找：布尔演算当时还缺乏转写大多数数学推理的手段，\(^{20}\) 因而对于莱布尼茨的伟大梦想只给出了一个很不完全的回答。建造一种更适合于数学的形式体系——其中变量和量词的引入（各自独立地归功于弗雷格 {[}117 a, b, c{]} 和 C. S. 皮尔士 {[}248 b{]}）构成主要的阶段——是逻辑学家和数学家的工作，与上述诸人不同，他们的首要目标恰恰是数学基础方面的应用。

弗雷格的计划 {[}117 b 与 c{]} 是要为算术建立一个基础，它基于一种借助“概念文字”（Begriffschrift）而形式化的逻辑；关于他定义自然数的方式，我们将在后面（第~29~页）再回来讨论。他的工作的特点，是对概念的分析极其精确、对细节极为讲究；正是由于这种倾向，他引入了许多后来在现代逻辑中显得极为重要的区分：例如，正是他第一个区分了命题的陈述与断言该命题为真，区分了属于关系与包含关系，区分了对象 x 与只含这个对象的集合 \(\{x\}\)，等等。他的形式化逻辑不仅包含数学中习用意义上的“变量”，还包含代表不确定关系的、可施以量词的“命题变元”，它后来（通过罗素和怀特海的工作）为元数学提供了基本工具。可惜他所采用的符号缺乏启发性，排版复杂得吓人，又远离数学家的实践；这就使数学家们望而却步，并大大削弱了弗雷格对其同时代人的影响。

皮亚诺的目标同时要宏大得多，也脚踏实地得多：出版一部完全用形式语言写成的“数学公式集”，其中不仅包含数理逻辑，而且包含数学各最重要分支的全部成果。他在一大群热心合作者（瓦伊拉蒂、皮耶里、帕多亚、瓦卡、维万蒂、法诺、布拉利-福尔蒂）的帮助下，以那样的速度完成了这一雄心勃勃的计划，本身就见证了他所采用符号体系的卓越：他的语言紧贴数学家的现行实践，引入了大量精心挑选的缩写符号，又特别由于一套用句点分隔符代替括号的巧妙办法，仍然相当易于阅读 {[}246 f{]}。皮亚诺创制的许多记号今天已被大多数数学家采用：我们举出 \(\in, \supset\)（但与当今的用法相反，其含义是“包含于”或“蕴涵”\(^{21}\)）、\(\cup, \cap, A - B\)（由差 \(a - b\) 组成的集合，其中 \(a \in A\)，\(b \in B\)）。另一方面，正是在《数学公式集》中，第一次出现了对函数一般概念以及直接像\(^{22}\) 与逆像概念的透彻分析，还有“序列不过是定义在 N 上的函数”这一评注。但是，量词在皮亚诺那里受到掣肘的限制（在他的系统中，原则上只能对形如 \(A \to B\)、\(A \Leftrightarrow B\) 或 \(A = B\) 的关系加量词）。此外，他的一些门徒近乎狂热的热情使他们大遭讥笑；尤其是 H. 庞加莱的批评（常常并不公正）给皮亚诺学派以沉重的打击，成为其学说在数学界传播的障碍。

有了弗雷格和皮亚诺，今天所使用的形式语言的基本要素便已齐备。流传最广的无疑是在罗素和怀特海的巨著《Principia Mathematica》中锤炼出来的那一种，它幸运地把弗雷格的精确与皮亚诺的便利结合在一起 {[}266{]}。现今大多数形式语言与它的区别，仅在于一些旨在简化其用法的次要改动。其中最巧妙的，我们举出关系的“函数式”书写（例如以 \(\in xy\) 代替 \(x \in y\)，这是武卡谢维奇想出来的，靠它括号可以完全省略）；但最有趣的无疑是希尔伯特对符号 \(\tau\) 的引入，它使人可以把量词 \(\exists\) 和 \(\forall\) 看作缩写记号，避免引入皮亚诺和罗素的“泛函”符号 \(\iota\)（它只适用于函数关系），最后还使人免去在集合论中陈述选择公理的必要（{[}163 a{]}，第 III 卷，第~183~页）。

